% !TeX program = xelatex
% !BIB program = biber
\documentclass[12pt]{article}
\usepackage[margin=1in]{geometry}

% Fuente Times New Roman (compilar con XeLaTeX)
\usepackage{fontspec}
\setmainfont{Times New Roman}

% Idioma e interlineado
\usepackage[spanish,es-nodecimaldot]{babel}
\usepackage{setspace}
\setstretch{1.15}

% Paquetes principales
\usepackage{amsmath,amssymb}
\usepackage{tikz-cd}
\usepackage{multicol}
\usepackage{graphicx}
\usepackage{booktabs}
\usepackage{caption}
\usepackage{titlesec}
\usepackage{enumitem}
\usepackage{longtable}
\usepackage{array}
\usepackage{xurl}
\usepackage{etoolbox}
\usepackage{csquotes}
\usepackage[
  backend=biber,
  style=apa,
  sorting=nyt,
  sortcites=true
]{biblatex}
\DeclareLanguageMapping{spanish}{spanish-apa}
\addbibresource{referencias.bib}

% Formato APA 7 para la lista de referencias.
\setlength{\bibhang}{0.5in}
\AtBeginBibliography{\setstretch{2}}

% Compatibilidad con longtable de TeX Live 2025: evita diagnósticos espurios
% de contracción infinita al dividir tablas entre varias páginas.
\makeatletter
\patchcmd{\LT@output}{\copy\LT@foot\vss}{\copy\LT@foot\vfil}{}{}
\patchcmd{\LT@output}{\copy\LT@foot\vss}{\copy\LT@foot\vfil}{}{}
\makeatother

% Títulos, leyendas y numeración en 12 puntos
\titleformat{\section}
  {\normalfont\normalsize\bfseries}{\thesection}{1em}{}
\titleformat{\subsection}
  {\normalfont\normalsize\bfseries}{\thesubsection}{1em}{}
\titleformat{\subsubsection}
  {\normalfont\normalsize\bfseries}{\thesubsubsection}{1em}{}
\captionsetup{font=normalsize,labelfont=bf}

% Párrafos
\setlength{\parindent}{1.5em}
\setlength{\parskip}{0pt}

% Numeración de páginas
\pagestyle{plain}

% Datos de la portada
\newcommand{\portadaUniversidad}{Universidad del Rosario}
\newcommand{\portadaPrograma}{Matemáticas aplicadas y ciencias de la computación}
\newcommand{\portadaAsignatura}{Análisis estadístico de datos}
\newcommand{\portadaTitulo}{Determinantes socioeconómicos e institucionales del desempeño en las pruebas Saber 11 (2020-2)}
\newcommand{\portadaIntegrantes}{%
  David Libreros\\
  Mateo Puente \\
  Nicolás González\\
  Santiago Díaz}
\newcommand{\portadaProfesora}{Luz Adriana Pineda}
\newcommand{\portadaFecha}{\today}
\newcommand{\portadaSemestre}{Segundo semestre académico}

\begin{document}

% Portada
\begin{titlepage}
  \thispagestyle{empty}
  \centering

  {\bfseries \portadaUniversidad\par}
  \vspace{0.6cm}
  {\portadaPrograma\par}
  \vspace{0.4cm}
  {\portadaAsignatura\par}

  \vfill

  {\bfseries \portadaTitulo\par}

  \vfill

  {
    \textbf{Integrantes}\par
    \vspace{0.3cm}
    \portadaIntegrantes\par
    \vspace{0.8cm}
    \textbf{Profesora}\par
    \vspace{0.3cm}
    \portadaProfesora\par
  }

  \vfill

  {
    \portadaSemestre\par
    \portadaFecha\par
  }
\end{titlepage}

% La numeración visible comienza después de la portada.
\setcounter{page}{1}

\section{Introducción}
El desempeño académico de los estudiantes de último grado de bachillerato en las pruebas Saber 11 en Colombia ha sido históricamente objeto de análisis por parte de instituciones educativas y entidades gubernamentales, debido al papel que puede ejercer como indicador de calidad educativa de la ciudad, la familia y la institución de donde proviene el estudiante, y de paso como insumo para el acceso a la educación superior por medio de incentivos del Estado o como indicador del nivel de enseñanza. Varios estudios respecto a este tema han señalado que el rendimiento en estas pruebas no depende únicamente de factores individuales del estudiante, sino que está mediado por el contexto socioeconómico y familiar, las características no solo académicas sino también sociales del establecimiento educativo y la ubicación geográfica en la que este se encuentra \parencite{garciagil2021}.

Este proyecto busca analizar los resultados de las pruebas Saber 11 del segundo semestre del 2020 (ICFES), por medio de una base de datos oficial extraída del portal de datos abiertos del Gobierno Nacional, con información sociodemográfica, familiar, institucional y de desempeño de aproximadamente 505.000 estudiantes de todo el territorio nacional. Resulta relevante estudiar este periodo porque corresponde al primer año en que la aplicación de las pruebas se realizó bajo condiciones de pandemia por el COVID-19, momento en el que las brechas relacionadas con los recursos tecnológicos y las condiciones del hogar (conectividad, materiales de estudio, acompañamiento familiar y académico) pudieron hacer más notorias las diferencias en el desempeño académico entre estudiantes de distintos contextos \parencite{cruz2016}.

Comprender cómo se asocian estas condiciones con el desempeño académico permite aportar evidencia para el diseño de políticas de equidad educativa y para la distribución de recursos hacia poblaciones en mayores condiciones de vulnerabilidad.

\section{Planteamiento del problema}
A pesar de que las pruebas Saber 11 buscan medir de manera estandarizada las competencias académicas de todos los estudiantes del país al finalizar el bachillerato, existe una problemática recurrente en la educación colombiana: los resultados no son homogéneos entre distintos grupos poblacionales. Los estudiantes de colegios oficiales frente a los de colegios privados, los de zonas rurales frente a los de zonas urbanas, o los provenientes de hogares con menor nivel socioeconómico, tienden a mostrar patrones marcados de diferencia en el desempeño \parencite{morales2024}.

Sin embargo, gran parte de los análisis realizados sobre estos resultados se limita a comparaciones descriptivas simples --promedios por departamento o por tipo de colegio--, sin explorar de manera conjunta los múltiples factores socioeconómicos, geográficos e institucionales que inciden en el rendimiento de cada estudiante. Esto deja un vacío en la comprensión de qué variables tienen mayor peso relativo y de si su efecto se mantiene al considerarlas simultáneamente en lugar de forma aislada.

Por tanto, este proyecto busca abordar dicho vacío explorando las conexiones y contrastes entre los individuos que presentan estas pruebas y su entorno, aplicando técnicas de estadística multivariada sobre la base de datos Saber 11 2020-2, con el fin de identificar los factores asociados a resultados poco equitativos en la población colombiana.

\section{Pregunta de investigación}
¿En qué medida el conjunto de características individuales, familiares, socioeconómicas, institucionales y geográficas registradas en la base de datos oficial Saber 11 2020-2 se asocia con el desempeño académico (puntaje global y por componente) de los estudiantes evaluados en Colombia?

\section{Objetivos}

\subsection{Objetivo general}
Determinar el grado de asociación entre las características individuales, familiares, institucionales y geográficas de los estudiantes y sus colegios, y el desempeño académico obtenido por los estudiantes de último año de bachillerato en las pruebas Saber 11 del segundo semestre de 2020, mediante técnicas de estadística multivariada.

\subsection{Objetivos específicos}
\begin{itemize}[leftmargin=1.8em, itemsep=0.4em, topsep=0.3em]
    \item Caracterizar demográfica e individualmente a la población evaluada (género, edad, pertenencia étnica, hábitos de lectura e internet, trabajo remunerado) y describir su posible relación con el desempeño obtenido.

    \item Analizar la asociación entre las condiciones familiares y socioeconómicas del hogar (estrato, composición del hogar, educación y ocupación de los padres, acceso a bienes y servicios, alimentación, índice socioeconómico individual) y el puntaje global y por componente.

    \item Explorar la relación entre las características institucionales y geográficas del colegio (naturaleza oficial o no oficial, calendario, jornada, bilingüismo, género del colegio, índice socioeconómico del establecimiento, zona urbana o rural y departamento de ubicación) y el desempeño de sus estudiantes.

    \item Identificar de manera conjunta cuáles de las variables anteriores tienen mayor peso relativo en la explicación del desempeño académico.
\end{itemize}

\section{Revisión bibliográfica}

El rendimiento académico de los estudiantes es un fenómeno que puede estar relacionado con diferentes características de su contexto individual, familiar y socioeconómico. En este sentido, la literatura previa permite identificar diversos factores que pueden estar asociados con el desempeño académico y que, por tanto, resultan relevantes para el análisis de los resultados obtenidos en las pruebas Saber 11.

En primer lugar, el estudio presentado en la \textit{Revista INFAD de Psicología} analiza la relación existente entre las variables socioeconómicas familiares y el rendimiento académico de estudiantes de educación secundaria obligatoria y bachillerato. Esta investigación resulta relevante para el presente proyecto debido a que permite establecer un antecedente sobre la importancia de las condiciones socioeconómicas del entorno familiar en el desempeño académico de los estudiantes \parencite{garciagil2021}. A partir de este antecedente, resulta pertinente considerar dentro del análisis variables relacionadas con las condiciones familiares y socioeconómicas de los estudiantes evaluados en las pruebas Saber 11 2020-2.

De manera más específica, \textcite{cruz2016} abordan la construcción de un índice socioeconómico familiar para los estudiantes que presentan la prueba Saber 11. Este trabajo constituye un antecedente especialmente relevante para el presente proyecto, debido a que utiliza información relacionada directamente con la población evaluada mediante las pruebas Saber 11 y estudia la relación entre las condiciones socioeconómicas de los estudiantes y su desempeño académico. La utilización de un índice socioeconómico permite, además, considerar que diferentes características del entorno familiar pueden estar relacionadas entre sí y que su análisis conjunto puede proporcionar una descripción más completa de las condiciones socioeconómicas de los estudiantes.

Este antecedente resulta particularmente útil para la investigación propuesta, ya que la base de datos Saber 11 2020-2 contiene información relacionada con las características individuales y familiares de los estudiantes, así como variables asociadas con sus condiciones socioeconómicas. Por esta razón, la información disponible permite explorar si las diferencias en estas características se encuentran asociadas con las diferencias observadas en los resultados obtenidos en la prueba.

Por otra parte, \textcite{morales2024} estudian la influencia del nivel socioeconómico de las familias sobre el rendimiento académico de los estudiantes. Este antecedente permite reforzar la importancia de analizar las condiciones socioeconómicas como una posible dimensión asociada con el desempeño académico. Asimismo, el estudio resulta de interés para el presente proyecto porque permite considerar que el rendimiento académico no necesariamente debe analizarse de manera aislada, sino teniendo en cuenta las condiciones del entorno en el que se desarrolla el estudiante.

La revisión de estos antecedentes permite identificar una relación común entre las condiciones socioeconómicas y familiares y el rendimiento académico. Sin embargo, estos estudios no abarcan necesariamente todas las dimensiones disponibles en la base de datos seleccionada para el presente proyecto. En particular, la base Saber 11 2020-2 permite incorporar simultáneamente características individuales, familiares, socioeconómicas, institucionales y geográficas, lo que posibilita ampliar el análisis hacia una perspectiva multivariada.

La fuente de datos seleccionada corresponde al conjunto oficial de resultados anonimizados de las pruebas Saber 11 de Calendario A del año 2020, publicado por el \textcite{icfes2026} a través del portal de datos abiertos del Gobierno de Colombia. La disponibilidad de información sobre diferentes características de los estudiantes y de los establecimientos educativos permite analizar el desempeño académico desde múltiples dimensiones.

En particular, esta base de datos resulta pertinente para el proyecto porque permite relacionar las características individuales y familiares de los estudiantes con los resultados obtenidos en las pruebas. De igual manera, la información disponible permite incorporar características relacionadas con el establecimiento educativo y con la ubicación geográfica, ampliando el análisis más allá de las condiciones exclusivamente individuales del estudiante.

En consecuencia, los antecedentes revisados proporcionan fundamentos para considerar que las características socioeconómicas y familiares constituyen dimensiones relevantes para el estudio del rendimiento académico. A su vez, el trabajo relacionado directamente con las pruebas Saber 11 evidencia la utilidad de este tipo de información para estudiar la relación entre las condiciones socioeconómicas y los resultados académicos. La base de datos Saber 11 2020-2 permite ampliar este enfoque al incorporar de manera conjunta diferentes dimensiones del contexto de los estudiantes y de las instituciones educativas.

Por lo anterior, la base de datos seleccionada presenta una alta pertinencia para la investigación propuesta, ya que permite estudiar la asociación entre un conjunto amplio de características y el desempeño académico de los estudiantes evaluados. Esto proporciona una base adecuada para la aplicación posterior de técnicas de estadística multivariada orientadas a identificar patrones, relaciones y variables que presenten una mayor asociación con el puntaje global y los resultados obtenidos en los diferentes componentes de las pruebas Saber 11.

\section{Marco metodológico}

\subsection{Fuente de los datos}

Los datos provienen de la base de datos pública \textit{Saber 11-2020-2} publicada por el Instituto Colombiano para la Evaluación de la Educación (ICFES) a través del portal de Datos Abiertos del Gobierno Nacional de Colombia (Instituto Colombiano para la Evaluación de la Educación - ICFES, Bogotá D.C., 2026). La base corresponde a los resultados consolidados de la aplicación de la prueba Saber 11 durante el segundo semestre calendario del año 2020, e incorpora información sociodemográfica del estudiante, características del hogar y del entorno familiar, atributos del establecimiento educativo donde culminó el bachillerato y, finalmente, los puntajes y percentiles obtenidos por el estudiante en cada una de las cinco pruebas que componen el examen, junto con el puntaje global y el percentil global. 

\subsection{Población y muestra}

La \textbf{población de referencia} corresponde a todos los estudiantes de último año de bachillerato del sistema educativo colombiano que presentaron la prueba Saber 11 durante el segundo semestre de 2020. La \textbf{población estudiada} está constituida por la totalidad de registros efectivamente publicados por el ICFES para esa cohorte. Dado que la prueba Saber 11 es de carácter obligatorio para todos los estudiantes de grado undécimo en el país y que la base de datos del ICFES reúne los resultados de todos los aspirantes (obligatorios y voluntarios) sin selección previa, la unidad de análisis no corresponde a una muestra probabilística sino a un \textbf{censo} de la población examinada en ese período.

\subsection{Número de observaciones}
El número total de registros contenidos en la base original corresponde a \textbf{504.872 observaciones}.

\subsection{Número de variables}
La base de datos original contiene un total de \textbf{81} variables por registro, de las cuales, tras el proceso de selección descrito en la sección~7, se conservaron \textbf{41} para el análisis estadístico multivariado: 40 variables de estudio y una empleada exclusivamente como criterio de filtro. A ellas se suma la variable derivada \texttt{EDAD}, construida durante la preparación de los datos.

\subsection{Contexto de recolección de los datos}
La información fue recolectada por el ICFES en el marco del proceso de aplicación de la prueba Saber 11 correspondiente al segundo semestre del año 2020. El contexto de pandemia por COVID-19 que caracterizó ese período le otorga a la cohorte analizada un interés particular: la aplicación de la prueba se realizó bajo protocolos sanitarios especiales y bajo condiciones excepcionales de virtualidad, conectividad y acompañamiento familiar en los hogares colombianos, por lo que las variables socioeconómicas --acceso a internet, computador, composición del hogar, dedicación a la lectura-- pueden mostrar una mayor variabilidad explicativa del desempeño que en cohortes prepandémicas. Adicionalmente, la base de datos corresponde a registros administrativos consolidados por el ICFES a partir de los formularios de inscripción y de las respuestas del examen, y no a una encuesta diseñada específicamente con fines de investigación; por tanto, su uso se enmarca en análisis secundarios de fuentes oficiales de carácter público. La manipulación y el análisis de los datos se realizan íntegramente en \texttt{R} mediante \textit{Quarto}, lo cual garantiza la trazabilidad y reproducibilidad del estudio.

\section{Descripción de las variables}

La base de datos Saber 11--2020-2 contiene un número elevado de columnas, no todas igualmente relevantes para responder la pregunta de investigación planteada. Por esta razón, la selección de variables se realizó de manera manual, conservando únicamente aquellas que aportan información sobre las características individuales del estudiante, su contexto familiar y socioeconómico, las particularidades del establecimiento educativo en el que cursa el último año de bachillerato y su desempeño académico. En la Tabla~\ref{tab:variables} se describen las columnas conservadas, agrupadas por categoría. Las columnas descartadas, junto con la justificación de su exclusión, se listan a continuación.

\subsection{Columnas descartadas}

De las 81 variables originales se descartaron las siguientes, agrupadas
según el criterio que motivó su exclusión:

\begin{itemize}[leftmargin=1.8em, itemsep=0.4em, topsep=0.3em]
  \item \textbf{Varianza nula o cardinalidad extrema:} \texttt{ESTU\_ESTUDIANTE},
  \texttt{ESTU\_PAIS\_RESIDE}, \texttt{PERIODO}, \texttt{ESTU\_PRIVADO\_LIBERTAD}
  y \texttt{COLE\_SEDE\_PRINCIPAL} presentan una única categoría o una
  distribución tan concentrada que no aportan información discriminante.
  Se descarta también \texttt{ESTU\_NACIONALIDAD}, con 99{,}4\,\% de casos en
  una sola categoría y 45 categorías de frecuencia inferior al 0{,}01\,\%.

  \item \textbf{Redundancia sobre el puntaje:} las variables
  \texttt{PERCENTIL\_*} y \texttt{DESEMP\_*} son reexpresiones ordinales de los
  puntajes \texttt{PUNT\_*}. Su inclusión conjunta generaría un primer
  componente artificial que reflejaría la dependencia funcional entre
  ellas y no la estructura del desempeño. Se conservan únicamente los
  puntajes continuos.

  \item \textbf{Redundancia geográfica:} el municipio de residencia, el de
  presentación de la prueba y el de ubicación del colegio coinciden en la
  mayoría de los registros. Su cardinalidad, superior a mil categorías,
  resulta además inmanejable en técnicas factoriales sobre datos
  categóricos. Se conserva \texttt{COLE\_DEPTO\_UBICACION}, con 33 categorías,
  junto con \texttt{COLE\_AREA\_UBICACION}.

  \item \textbf{Faltante estructural:} \texttt{ESTU\_TIPOREMUNERACION} solo aplica
  a quienes declaran trabajo remunerado. La ausencia de respuesta no
  constituye un dato perdido sino una pregunta no aplicable.

  \item \textbf{Índice duplicado:} se descarta
  \texttt{ESTU\_NSE\_INDIVIDUAL} porque discretiza la variable continua
  \texttt{ESTU\_INSE\_INDIVIDUAL}, que se conserva.

  \item \textbf{Inadecuación al diseño:} \texttt{ESTU\_GENERACION-E} identifica a
  beneficiarios de un programa cuya asignación es posterior a la prueba y
  depende del puntaje obtenido.

  \item \textbf{Identificadores:} \url{ESTU_CONSECUTIVO},
  \url{COLE_CODIGO_ICFES}, \url{COLE_COD_DANE_ESTABLECIMIENTO},
  \url{COLE_COD_DANE_SEDE}, \url{COLE_NOMBRE_ESTABLECIMIENTO} y
  \url{COLE_NOMBRE_SEDE} identifican registros y establecimientos, pero no
  constituyen variables de estudio: los códigos carecen de contenido
  sustantivo y los nombres presentan variantes ortográficas para un mismo
  establecimiento. El identificador \url{ESTU_CONSECUTIVO} se emplea
  únicamente en la verificación de registros duplicados, sobre la base
  original.
\end{itemize}

\subsection{Descripción de las variables seleccionadas}

\newcolumntype{L}[1]{>{\raggedright\arraybackslash}p{#1}}
\begin{longtable}{L{2.2cm} L{6.5cm} L{6.3cm}}
  \caption{Descripción de las variables seleccionadas para el análisis.}\label{tab:variables}\\
  \toprule
  \textbf{Categoría} & \textbf{Variable} & \textbf{Descripción} \\
  \midrule
  \endfirsthead

  \midrule
  \multicolumn{3}{r}{\textit{Continúa en la página siguiente...}} \\
  \midrule
  \textbf{Categoría} & \textbf{Variable} & \textbf{Descripción} \\
  \midrule
  \endhead

  \midrule
  \multicolumn{3}{r}{\textit{Continúa en la página siguiente...}} \\
  \endfoot

  \bottomrule
  \endlastfoot

  % --- Características demográficas ---
  \midrule
  \multicolumn{3}{l}{\textit{Características demográficas}} \\
  ESTU & \texttt{ESTU\_GENERO} & Género reportado por el estudiante. \\
  ESTU & \texttt{ESTU\_FECHANACIMIENTO} & Fecha de nacimiento; se utiliza para calcular la edad al momento de presentar la prueba. \\
  ESTU & \texttt{ESTU\_TIENEETNIA} & Indica si el estudiante se autorreconoce como perteneciente a algún grupo étnico. \\

  % --- Hábitos y dedicación ---
  \midrule
  \multicolumn{3}{l}{\textit{Hábitos y dedicación}} \\
  ESTU & \texttt{ESTU\_DEDICACIONLECTURADIARIA} & Tiempo diario dedicado a la lectura fuera del ámbito escolar. \\
  ESTU & \texttt{ESTU\_DEDICACIONINTERNET} & Tiempo diario dedicado al uso de internet. \\
  ESTU & \texttt{ESTU\_HORASSEMANATRABAJA} & Número de horas semanales dedicadas a trabajo remunerado. \\

  % --- Composición y condiciones del hogar ---
  \midrule
  \multicolumn{3}{l}{\textit{Composición y condiciones del hogar}} \\
  FAMI & \texttt{FAMI\_ESTRATOVIVIENDA} & Estrato socioeconómico de la vivienda (1 a 6, según clasificación del DANE). \\
  FAMI & \texttt{FAMI\_PERSONASHOGAR} & Número total de personas que habitan el hogar. \\
  FAMI & \texttt{FAMI\_CUARTOSHOGAR} & Número de cuartos con que cuenta la vivienda. \\
  FAMI & \texttt{FAMI\_SITUACIONECONOMICA} & Percepción de la situación económica del hogar declarada por el estudiante. \\

  % --- Capital educativo de los padres ---
  \midrule
  \multicolumn{3}{l}{\textit{Capital educativo y ocupacional de los padres}} \\
  FAMI & \texttt{FAMI\_EDUCACIONPADRE} & Nivel educativo máximo alcanzado por el padre o figura paterna (ordinal). \\
  FAMI & \texttt{FAMI\_EDUCACIONMADRE} & Nivel educativo máximo alcanzado por la madre o figura materna (ordinal). \\
  FAMI & \texttt{FAMI\_TRABAJOLABORPADRE} & Tipo de ocupación del padre o figura paterna. \\
  FAMI & \texttt{FAMI\_TRABAJOLABORMADRE} & Tipo de ocupación de la madre o figura materna. \\
  FAMI & \texttt{FAMI\_NUMLIBROS} & Cantidad aproximada de libros presentes en el hogar (ordinal). \\

  % --- Tenencia de bienes y acceso tecnológico ---
  \midrule
  \multicolumn{3}{l}{\textit{Tenencia de bienes y acceso tecnológico}} \\
  FAMI & \texttt{FAMI\_TIENEINTERNET} & Indica si el hogar cuenta con servicio de internet. \\
  FAMI & \texttt{FAMI\_TIENECOMPUTADOR} & Indica si el hogar cuenta con computador de escritorio o portátil. \\
  FAMI & \texttt{FAMI\_TIENEAUTOMOVIL} & Indica si el hogar posee automóvil particular. \\
  FAMI & \texttt{FAMI\_TIENEMOTOCICLETA} & Indica si el hogar posee motocicleta. \\
  FAMI & \texttt{FAMI\_TIENELAVADORA} & Indica si el hogar posee lavadora. \\
  FAMI & \texttt{FAMI\_TIENEHORNOMICROOGAS} & Indica si el hogar posee horno microondas o estufa de gas. \\
  FAMI & \texttt{FAMI\_TIENECONSOLAVIDEOJUEGOS} & Indica si el hogar posee consola de videojuegos. \\

  % --- Seguridad alimentaria ---
  \midrule
  \multicolumn{3}{l}{\textit{Seguridad alimentaria del hogar}} \\
  FAMI & \texttt{FAMI\_COMELECHEDERIVADOS} & Frecuencia de consumo de leche y derivados en el hogar. \\
  FAMI & \texttt{FAMI\_COMECARNEPESCADOHUEVO} & Frecuencia de consumo de carne, pescado o huevo en el hogar. \\
  FAMI & \texttt{FAMI\_COMECEREALFRUTOSLEGUMBRE} & Frecuencia de consumo de cereales, frutos secos o legumbres en el hogar. \\

  % --- Establecimiento educativo ---
  \midrule
  \multicolumn{3}{l}{\textit{Características del establecimiento educativo}} \\
  COLE & \texttt{COLE\_NATURALEZA} & Naturaleza jurídica del colegio (oficial o no oficial). \\
  COLE & \texttt{COLE\_CALENDARIO} & Calendario académico bajo el cual opera el colegio (A o B). \\
  COLE & \texttt{COLE\_GENERO} & Tipo de colegio según el género de los estudiantes (masculino, femenino o mixto). \\
  COLE & \texttt{COLE\_BILINGUE} & Indica si el colegio se encuentra registrado como bilingüe ante el MEN. \\
  COLE & \texttt{COLE\_JORNADA} & Jornada en la que funciona el colegio (mañana, tarde, noche, única, sabatina). \\
  COLE & \texttt{COLE\_AREA\_UBICACION} & Área de ubicación del colegio (urbana o rural). \\
  COLE & \texttt{COLE\_DEPTO\_UBICACION} & Departamento donde se encuentra ubicado el colegio. \\

  % --- Índices socioeconómicos ---
  \midrule
  \multicolumn{3}{l}{\textit{Índices socioeconómicos}} \\
  ESTU & \texttt{ESTU\_INSE\_INDIVIDUAL} & Índice socioeconómico individual calculado por el ICFES (continuo). \\
  ESTU & \texttt{ESTU\_NSE\_ESTABLECIMIENTO} & Nivel socioeconómico promedio del establecimiento educativo (ordinal). \\

  % --- Desempeño académico ---
  \midrule
  \multicolumn{3}{l}{\textit{Desempeño académico (variables de respuesta)}} \\
  PUNT & \texttt{PUNT\_LECTURA\_CRITICA} & Puntaje obtenido en la prueba de lectura crítica (0--100). \\
  PUNT & \texttt{PUNT\_MATEMATICAS} & Puntaje obtenido en la prueba de matemáticas (0--100). \\
  PUNT & \texttt{PUNT\_C\_NATURALES} & Puntaje obtenido en la prueba de ciencias naturales (0--100). \\
  PUNT & \texttt{PUNT\_SOCIALES\_CIUDADANAS} & Puntaje obtenido en la prueba de sociales y ciudadanas (0--100). \\
  PUNT & \texttt{PUNT\_INGLES} & Puntaje obtenido en la prueba de inglés (0--100). \\
  PUNT & \texttt{PUNT\_GLOBAL} & Puntaje global del estudiante, calculado como promedio ponderado de las cinco pruebas. \\

  % --- Criterio de filtro ---
  \midrule
  \multicolumn{3}{l}{\textit{Criterio de filtro}} \\
  ESTU & \texttt{ESTU\_ESTADOINVESTIGACION} & Estado del registro dentro del proceso de investigación del ICFES; se conservan únicamente los registros publicables. \\
\end{longtable}

\section{Preparación de la base de datos}
Para garantizar la calidad de la información, se exploraron los datos con el fin de revisar registros duplicados, datos faltantes y valores inválidos. A continuación se describe el proceso seguido en \texttt{R} junto con algunas observaciones sobre las decisiones tomadas.

\subsection{Datos inválidos}
Se comenzó revisando los registros validados por el ICFES, es decir, aquellos que no estuvieran marcados como \texttt{VALIDEZ OFICINA JURÍDICA}, categoría que corresponde a registros que en su momento estuvieron bajo investigación. Se trata de un procedimiento de rigor más que de un filtro con efecto apreciable, dado el tamaño de la base: se encontraron 44 entradas marcadas de esta manera, las cuales se descartaron.

\subsection{Datos duplicados}
Utilizando la base completa, que incluye los identificadores, no se encontraron registros duplicados: el identificador \texttt{ESTU\_CONSECUTIVO} es único en las 504.872 observaciones originales.

\subsection{Datos faltantes}
La revisión de los datos faltantes comenzó por cuantificar su magnitud. La variable con mayor proporción de ausencias es \texttt{COLE\_BILINGUE}, con un 16,4\,\% de los registros (83.003), cifra considerable. Se tomó la decisión de no filtrar dichos registros, pues no se justifica afectar la totalidad de los análisis posteriores por una sola variable. Además, el valor ausente coexiste con las categorías \texttt{S} y \texttt{N}, de modo que la ausencia de información no implica que el colegio no sea bilingüe, sino que no se dispone del registro ante el Ministerio de Educación Nacional.

Por otra parte, 23 variables del cuestionario socioeconómico presentan entre un 2,6\,\% y un 4\,\% de datos faltantes. Al agrupar los registros según el número de preguntas del cuestionario sin responder se identificaron tres poblaciones diferenciadas: un 89,2\,\% con el cuestionario completo, un 8,3\,\% con ausencias puntuales --de las cuales cerca del 74\,\% corresponde a un único ítem sin responder-- y un 2,5\,\% que prácticamente no diligenció el cuestionario. Esta última proporción coincide con los registros sin valor en \texttt{ESTU\_INSE\_INDIVIDUAL}, lo cual es consistente con el hecho de que el ICFES calcula dicho índice a partir de esas mismas respuestas.

La comparación de las características observables de los tres grupos muestra que la ausencia de información no es aleatoria. El grupo con no respuesta parcial registra un puntaje global promedio de 228,8 puntos, frente a 250,4 del grupo completo, y presenta una mayor proporción de estudiantes de colegios oficiales (81,6\,\% frente a 77,3\,\%) y de zonas rurales (20,6\,\% frente a 16,1\,\%). El mecanismo de pérdida se clasifica, en consecuencia, como \textit{missing at random} (MAR) y no como completamente aleatorio.

Se conservan únicamente los registros sin valores faltantes en las 22 variables del cuestionario socioeconómico (n = 450.145; 89,2\,\%), dado que las técnicas multivariadas previstas requieren observaciones completas y de esta manera se garantiza que todos los análisis se apliquen sobre la misma población. Dado el sesgo documentado en el párrafo anterior, las brechas socioeconómicas que se estimen constituyen muy probablemente una cota inferior de las brechas reales.

\subsection{Recodificación}
A partir de la consulta de niveles realizada durante el análisis de datos faltantes, se recodificaron las columnas categóricas como factores, fijando explícitamente el orden de los niveles en las variables ordinales. Este paso resulta necesario para el manejo correcto de los datos en \texttt{R}: al declarar los niveles se evita que las funciones de modelamiento y graficación los ordenen alfabéticamente y se preserva la escala de las variables ordinales, condición indispensable para la interpretación de los ejes en las técnicas factoriales.

Cabe señalar que las categorías \texttt{NO SABE} y \texttt{NO APLICA} en las variables de educación de los padres, así como la categoría \texttt{SIN ESTRATO}, no se tratan como valores faltantes: corresponden a respuestas sustantivas --desconocimiento del nivel educativo, ausencia de la figura parental y vivienda no estratificada, respectivamente-- y se conservan como categorías propias fuera de la escala ordinal.

\subsection{Creación de la columna de edad}
Dado que se dispone de la fecha de nacimiento, se construyó la variable \texttt{EDAD} como la diferencia entre esa fecha y la de aplicación de la prueba. La distribución resultante es coherente con la población objetivo, con una mediana de 17,4 años y un rango intercuartílico de 16,9 a 18,2 años. Se detectaron 183 registros (0,041\,\%) con valores fuera del rango plausible, definido entre 10 y 100 años, atribuibles a errores en la fecha registrada --entre ellos un valor centinela de 1899 en diez casos--. Dichos valores se codificaron como faltantes y su tratamiento se deja al criterio del método utilizado en cada análisis posterior.

\section{Conclusiones preliminares}

Esta entrega corresponde a la fase de selección y preparación de la base de datos, por lo que las conclusiones se refieren a la calidad de la información disponible y a las condiciones bajo las cuales podrán aplicarse las técnicas multivariadas, y no aún a la relación entre las variables de estudio.

La base \textit{Saber 11--2020-2} presenta una calidad alta para un registro administrativo de gran escala: no se identificaron duplicados, únicamente 44 registros resultaron no publicables y los valores inválidos de fecha de nacimiento representan apenas un 0,041\,\% del total. Tras aplicar los criterios de validez y completitud, la base de trabajo quedó constituida por 450.145 registros y 41 variables, a las que se suma la variable derivada \texttt{EDAD}.

El principal hallazgo de esta fase corresponde a la estructura de la no respuesta. La ausencia de información en el cuestionario socioeconómico no es aleatoria, sino que se concentra en estudiantes de menor desempeño y de colegios oficiales y rurales. La exclusión de estos registros, que es necesaria para aplicar técnicas que requieren observaciones completas, implica un sesgo de selección que debe tenerse presente en la interpretación de los resultados, siendo estas que las brechas que se estimen constituyen una cota inferior de las existentes en la población.

La revisión de la distribución de categorías reveló, además, que tres variables del establecimiento educativo presentan un desbalance que restringe su uso: \texttt{COLE\_CALENDARIO} concentra un 0,07\,\% de los registros en la categoría B, \texttt{COLE\_GENERO} un 0,79\,\% en colegios masculinos y \texttt{COLE\_BILINGUE} un 1,03\,\% en colegios bilingües. Las categorías de frecuencia muy baja concentran una inercia desproporcionada en las técnicas factoriales y pueden generar ejes que reflejen el comportamiento de unos pocos casos en lugar de la estructura general de los datos, por lo que se prevé incorporarlas como variables suplementarias.

Finalmente, dado que la unidad de análisis corresponde a un censo y que el número de observaciones supera las 450.000, las pruebas de significancia estadística resultarán poco informativas: diferencias sustantivamente irrelevantes alcanzarán significancia. La interpretación de los resultados deberá apoyarse en medidas de tamaño de efecto y en la proporción de varianza explicada.

Con estas salvedades, la base seleccionada resulta adecuada para responder la pregunta planteada: el bloque de variables familiares ofrece una cobertura amplia de las dimensiones del capital del hogar, el bloque de desempeño permite un tratamiento multivariado de la variable respuesta en lugar de reducirla al puntaje global, y el índice socioeconómico calculado por el ICFES ofrece un criterio de contraste externo para los resultados obtenidos.

\section{Declaración de uso de inteligencia artificial}
Se usaron herramientas de Inteligencia Artificial (IA) para el desarrollo de esta entrega, para fines como mejorar la redacción, 
depurar código, documentación sobre el uso de herramientas como Quarto y R, y replantear enfoques o decisiones posiblemente erroneas.
En general se buscó con estas herramientas elevar la calidad del trabajo, manteniendo estructuralmente el criterio personal.

Entre las herramientas utilizadas se encuentran:
\begin{itemize}
  \item Claude (Opus 5)
  \item Prism
\end{itemize}


\printbibliography[title={Referencias}]

\end{document}